\section{Classification of rational bilinear tail certificates}
\label{sec:bilinear}

The latest alternating-tail report asks whether allowing both adjacent
products gives a useful classification for two different denominator
polynomials. We answer that question exactly. The classification is
relative to an explicit certificate class; failure of its criterion
does not exclude longer telescopers or other special-function identities.

\subsection{The certificate class and its factorization criterion}

Let $P,Q\in\C[z]$ be nonzero, let $\alpha,\beta,t\ne0$, and suppose
two sequences satisfy
\begin{equation}\label{bi:recurrences}
 E_{n-1}=E_n+\frac{\alpha^n}{P(n)},\qquad
 F_{n-1}=F_n+\frac{\beta^n}{Q(n)}.
\end{equation}
For rational functions $A,B$, consider
\begin{equation}\label{bi:certificate}
 T_n=t^n\{A(n+1)E_nF_{n+1}+B(n+1)E_{n+1}F_n\}.
\end{equation}
We ask that, after using \eqref{bi:recurrences}, the coefficients of
$E_n$ and $F_n$ in $T_n-T_{n-1}$ both vanish identically. The remaining
term is then a rational weight times $E_nF_n$. This is an algebraic
condition on the certificate, independent of any conjectured numerical
relation among the tails.

\begin{theorem}[Complete two-adjacent-product classification]
\label{bi:classification}
A nonzero rational certificate \eqref{bi:certificate} has the stated
cancellation property if and only if
\begin{equation}\label{bi:criterion}
 t^2\alpha\beta=1,\qquad
 \frac{P(z)}{Q(z)}=\rho h(z)h(z+1)
\end{equation}
for a nonzero constant $\rho$ and a nonzero rational function $h$.
In that case all certificates for the chosen $t$ are scalar multiples
of
\begin{equation}\label{bi:AB}
 A(z)=Q(z)h(z),\qquad B(z)=\kappa Q(z)h(z+1),\qquad
 \kappa=-t\beta.
\end{equation}
Writing $C=A+B$, their exact difference is
\begin{equation}\label{bi:difference}
 T_n-T_{n-1}
   =t^n\{C(n+1)-t^{-1}C(n)\}E_nF_n.
\end{equation}
Thus, when the rational factorization exists, there are exactly two
outer-color choices $t$, and the certificate space for either choice
is one-dimensional.
\end{theorem}

Here ``nonzero'' means that $(A,B)$ is not the zero pair. At a special
parameter value the resulting weight can vanish identically, or the
evaluated tails can make the certificate zero. Such degeneracies do not
change the formal cancellation classification. An example is the
equal-shift limit of the alternating kernel below, whose nontrivial
diagonal identity is recovered by holomorphic continuation.

\begin{proof}
Direct expansion gives the quadratic coefficient
$A(n+1)+B(n+1)-t^{-1}\{A(n)+B(n)\}$. The linear coefficients vanish
exactly when
\begin{equation}\label{bi:linear}
 \beta\frac{A(z+1)}{Q(z+1)}+t^{-1}\frac{B(z)}{Q(z)}=0,
 \qquad
 \alpha\frac{B(z+1)}{P(z+1)}+t^{-1}\frac{A(z)}{P(z)}=0.
\end{equation}
If one of $A,B$ is zero, both are zero. Otherwise set
$h=A/Q$ and $R=P/Q$. The first equation gives
$B=-t\beta Qh(z+1)$; substitution into the second gives
\begin{equation}\label{bi:twostep}
 \frac{R(z+1)}{R(z)}
       =t^2\alpha\beta\frac{h(z+2)}{h(z)}.
\end{equation}
A nonzero rational function $r$ satisfies $r(z+c)/r(z)\to1$ as
$z\to\infty$ for every fixed $c$. Therefore
$t^2\alpha\beta=1$. Equation \eqref{bi:twostep} now says that
$R(z)/\{h(z)h(z+1)\}$ is a rational function with period one.
A nonconstant rational periodic function would have infinitely many
poles, or, if polynomial, could not have a nonzero period. It is
therefore the constant $\rho$. This proves necessity.

Conversely, \eqref{bi:criterion} and \eqref{bi:AB} verify
\eqref{bi:linear} by substitution, proving the difference identity.
If $h_1,h_2$ are two solutions, their quotient $r$ satisfies
$r(z)r(z+1)=\text{constant}$ and hence $r(z+2)=r(z)$.
It is constant by the same rational-periodicity argument. This proves
the uniqueness assertion.
\end{proof}

\subsection{An effective criterion on the divisor of the rational function}

The condition in \eqref{bi:criterion} does not require an infinite
search for $h$. Reduce $R=P/Q$ first, and group its finite zeros and
poles into equivalence classes modulo integer translation. For an orbit
$\xi+\Z$, put
\[
 r_j=\ord_{z=\xi+j}R(z).
\]
Only finitely many $r_j$ are nonzero; poles have negative multiplicity.

\begin{proposition}[Alternating divisor test]
\label{bi:divisor}
The factorization $R(z)=\rho h(z)h(z+1)$ with rational $h$ exists
if and only if, for every such orbit,
\begin{equation}\label{bi:divisorcondition}
 \sum_{j\in\Z}(-1)^j r_j=0.
\end{equation}
When this holds, the multiplicity of $h$ at $\xi+j$ is
\begin{equation}\label{bi:divisorconstruction}
 h_j=\sum_{k\ge0}(-1)^k r_{j+k}.
\end{equation}
The constant $\rho$ then matches the leading coefficients at infinity.
\end{proposition}

\begin{proof}
The multiplicities must satisfy $r_j=h_j+h_{j+1}$.
The unique solution vanishing above the support of $r$ is
\eqref{bi:divisorconstruction}. It also vanishes below that support
exactly when \eqref{bi:divisorcondition} holds. This constructs a
rational $h$ with the required finite divisor. The quotient
$R/(hh(\cdot+1))$ has neither finite zeros nor poles, and its order
at infinity is also zero, because $\sum r_j=2\sum h_j$.
It is therefore a nonzero constant. Changing the orbit representative
only multiplies \eqref{bi:divisorcondition} by a sign.
\end{proof}

This is a finite algebraic test on the reduced rational function. In
particular, its degree at infinity must be even. Generic quadratic
polynomials with four zeros in distinct integer-translation orbits
fail the test. By contrast,
\begin{equation}\label{bi:adjacent-ratio}
 \frac{(z+a)(z+a+1)}{(z+b)(z+b+1)}
       =\frac{z+a}{z+b}\frac{z+a+1}{z+b+1}
\end{equation}
passes for arbitrary independent $a,b$. This example gives the
two-shift identities in the next section. The divisor argument is an
elementary instance of rational difference-equation reasoning; the
contribution here is its exact use to classify this bilinear tail
certificate and all compatible colors.

\subsection{The finite endpoint and the infinite identity}

Assume now that $\deg P=d\ge2$, $\deg Q=e\ge2$, and neither polynomial
vanishes at a positive integer. Let
\[
 0<|\alpha|,|\beta|\le1,\qquad \alpha\ne1,\quad\beta\ne1,
\]
and define the ordinary tails
\begin{equation}\label{bi:tails}
 E_n=\sum_{k>n}\frac{\alpha^k}{P(k)},\qquad
 F_n=\sum_{k>n}\frac{\beta^k}{Q(k)}.
\end{equation}
Suppose the equivalent conditions of Theorem~\ref{bi:classification}
hold. A zero or pole of $h$ at a positive integer would force, by
\eqref{bi:criterion}, alternating nonzero multiplicities at all later
integers, which is impossible for a rational function. Consequently
$h(1)$ is finite and nonzero. The rational functions $A,B$ are regular
at nonnegative integers after removal of any removable singularities.

\begin{theorem}[Bilinear summation with its exact upper endpoint]
\label{bi:sum}
Put $W(n)=C(n+1)-t^{-1}C(n)$ and
\begin{equation}\label{bi:endpoint}
 \mathcal E=
 \beta h(1)E_0-\frac{t^{-1}}{\rho h(1)}F_0
                  -t^{-1}C(0)E_0F_0.
\end{equation}
For every integer $N\ge0$,
\begin{equation}\label{bi:finite}
 \sum_{n=0}^N t^n W(n)E_nF_n=\mathcal E+T_N.
\end{equation}
Moreover,
\begin{equation}\label{bi:infinite}
 \boxed{\displaystyle
 \sum_{n=0}^{\infty}t^n W(n)E_nF_n=\mathcal E,}
\end{equation}
and the series converges absolutely. For fixed $P,Q,h$, it converges
locally normally in the admissible colors away from $1$, with either
locally chosen branch of $t$.
\end{theorem}

\begin{proof}
Sum \eqref{bi:difference} for $1\le n\le N$. Since
$E_1=E_0-\alpha/P(1)$ and $F_1=F_0-\beta/Q(1)$,
\[
 T_0=C(1)E_0F_0-\beta h(1)E_0
                         -\frac{\alpha\kappa}{\rho h(1)}F_0.
\]
Use $\alpha\kappa=-t^{-1}$ and restore the $n=0$ summand.
This gives \eqref{bi:finite}, including the lower-endpoint term
$C(0)E_0F_0$.

Summation by parts, using bounded partial geometric sums, gives
$E_n=O(|\alpha|^{n+1}n^{-d})$ and
$F_n=O(|\beta|^{n+1}n^{-e})$. The bounds are locally uniform as
stated. The factorization forces $\deg h=(d-e)/2$, so $A,B,C$
are $O(n^{(d+e)/2})$. Since $|t|=|\alpha\beta|^{-1/2}$,
the summand and $T_N$ are respectively bounded by a constant times
\[
 |\alpha\beta|^{n/2}n^{-(d+e)/2},\qquad
 |\alpha\beta|^{N/2}N^{-(d+e)/2}.
\]
The first is summable and the second tends to zero. This proves the
infinite identity and the local normal convergence assertion.
\end{proof}

The finite identity remains true outside these sufficient convergence
hypotheses whenever the recurrences make sense. Any additional ordinary
or regularized limit needs its own endpoint analysis.

\subsection{A Hermitian polylogarithm-tail law}

The two inner colors need not agree. Taking conjugate colors gives an
Hermitian identity with explicit dependence on the color.  Its
weights are positive for the power-polynomial specialization below.

\begin{corollary}[Conjugate-color completion]
\label{bi:hermitian}
Let $P$ be a real polynomial of degree at least two with no positive
integer zero. For $0<\theta<2\pi$, write
\[
 E_n(\theta)=\sum_{k>n}\frac{e^{ik\theta}}{P(k)}.
\]
Then
\begin{equation}\label{bi:hermitian-formula}
 \sum_{n\ge0}\{P(n+1)-P(n)\}|E_n(\theta)|^2
 =\cot\frac\theta2\,\Im E_0(\theta)-\Re E_0(\theta)
                         -P(0)|E_0(\theta)|^2.
\end{equation}
For every integer $d\ge2$, in particular,
\begin{equation}\label{bi:polylog-hermitian}
 \boxed{\begin{aligned}
 &\sum_{n\ge0}\{(n+1)^d-n^d\}
 \left|\Li_d(e^{i\theta})-\sum_{k=1}^n\frac{e^{ik\theta}}{k^d}\right|^2\\
 &\hspace{18mm}=\cot\frac\theta2\,\Im\Li_d(e^{i\theta})
                                  -\Re\Li_d(e^{i\theta}).
 \end{aligned}}
\end{equation}
\end{corollary}

\begin{proof}
In Theorem~\ref{bi:sum}, take $Q=P$, $h=1$, $\rho=1$,
$\alpha=e^{i\theta}$, $\beta=e^{-i\theta}$, and $t=1$.
Then $C=(1-e^{-i\theta})P$ and $F_n=\overline{E_n}$.
Divide by $1-e^{-i\theta}$ and use
\[
 \frac{e^{-i\theta}E-\overline E}{1-e^{-i\theta}}
       =-\Re E+\cot(\theta/2)\Im E.
\]
\end{proof}

At the Gaussian color and order two, this gives the explicit identity
\begin{equation}\label{bi:gaussian-example}
 \boxed{\displaystyle
 \sum_{n\ge0}(2n+1)
 \left|\Li_2(i)-\sum_{k=1}^n\frac{i^k}{k^2}\right|^2
          =G+\frac{\pi^2}{48},}
\end{equation}
where $G=\sum_{k\ge0}(-1)^k/(2k+1)^2$ is Catalan's constant.
Here $\Li_2(i)=-\pi^2/48+iG$, obtained by separating even and odd
terms of the defining series. At $\theta=\pi$ the formula recovers
the ordinary-outer alternating square law of \cite{tail-alternating}.
For $d>2$, its limit as $\theta\to0$ is also an ordinary identity:
\[
 \sum_{n\ge0}\{(n+1)^d-n^d\}
                 \{\zeta(d)-H_n^{(d)}\}^2
            =2\zeta(d-1)-\zeta(d).
\]
Indeed the absolute tail is bounded by $O(n^{1-d})$, giving a
summable majorant $O(n^{1-d})$ after multiplication by the weight.
The right side follows by differentiating the color at zero, justified
by $d>2$. This limit is not asserted at $d=2$, where the ordinary
series diverges. The new classification supplies these formulas as
corollaries; they are not separate claims of global literature priority.
